\documentclass{article}
\usepackage{graphicx} % Required for inserting images
\usepackage{datetime} % Add this line to include the datetime package
\usepackage{amsmath}
\usepackage{amssymb}
\usepackage{amsthm}
\usepackage{tikz}
%\usepackage{hyperref}
\usepackage{xcolor} % Required for defining custom colors

\definecolor{darkblue}{rgb}{0.0, 0.0, 0.5} % Define a darker blue color

\usepackage[colorlinks=true, linkcolor=darkblue, urlcolor=darkblue, citecolor=darkblue]{hyperref} % Use the darker blue color

% Define theorem style
\theoremstyle{plain} % Bold title, italic body
\newtheorem{theorem}{Theorem}
\newtheorem{lemma}[theorem]{Lemma}
\newtheorem{corollary}[theorem]{Corollary}

\theoremstyle{definition} % Bold title, normal body
\newtheorem{definition}[theorem]{Definition}
\newtheorem{example}[theorem]{Example}


\title{Homework 3}
\author{Aaron Honculada}
\date{\monthname[\month] \the\year}

\begin{document}

\maketitle

\section*{Problem 1}
\noindent The \textit{cycle type} of a permutation $\sigma \in S_n$ is the sequence of the lengths of
its cycles written
in weakly decreasing order. For instance, the permuation $(1\text{ }2)(3\text{ }4\text{ }5\text{ }6)(8\text{ }9) \in S_9$ has
cycle type $(4,2,2,1)$.


\noindent\textit{Show that for all $\sigma, \tau \in S_n$, the permutations $\sigma\text{ and }\tau^{-1}\sigma\tau$ have
the same cycle type.}

\begin{proof}
    We know that every permutation $\sigma$ can be expressed as the product of disjoint cycles $\sigma_1, \sigma_2, \cdots, \sigma_k$.
    It also has been shown that conjugation is a group homomorphism, where $\tau^{-1}(\sigma_i\sigma_j)\tau = \tau^{-1}\sigma_i\tau \cdot
    \tau^{-1}\sigma_j\tau$. If permutation $\sigma$ has a cycle type of $n$, then $n$ can be partitioned into the cycle types of the disjoint
    cycles whose products are $\sigma$. Otherwise stated, if $\sigma=\sigma_1\sigma_2\cdots\sigma_k$ and $\sigma \in S_n$, then cycle
    $\sigma_i$ has cycle type $n_i$ and $n = \sum_{i=1}^{k}n_i$. Since each $\sigma_i,\sigma_j$ is disjoint,
    then the cycle type of $\sigma_{ij} = \sigma_i\sigma_j$ is $(n_i, n_j)$. Therefore the cycle type of $\sigma = 
    \tau^{-1}\sigma\tau$.

    
\end{proof}



\section*{Problem 2}
\noindent\textit{Let $n \geq 2$ be an integer. Show that the symmetric group $S_n$ is generated by the two
elements $(1 \text{ }2)$ and $(1\text{ }2\text{ }3\cdots n)$.}

\begin{proof}
    When $n=2$, we see that all of $S_2$ can be generated by $ \langle (1\text{ }2)\rangle =\{e, (1\text{ }2)\}$. Now we consider when 
    $n > 2$.

    Let $G= \langle (1 \text{ }2), (1\text{ }2 \cdots n) \rangle$ and $s=(1\text{ }2\cdots n)$. Each time $s$ is applied it rotates the elements by one position. Given $s$ is a $n$-cycle, then applying $s$
    $n-1$ times will reverse the order of the elements and produce an inverse of $s$. Therefore $s^{n-1}=s^{-1}$ and $s^{-1} \in G$. 
    
    By conjugating $(1\text{ }2)$ by powers of $s$ where we have $s^{-k}(1\text{ }2)s^k$, 
    we are able to generate all transpositions of the form $(i, i+1)$ where $1 \leq i < n$.
    
    Thus all transpositions of adjacent integers from 1 to $n$ are in $\langle (1\text{ }2), s \rangle$.
    
    For $k > 2$, we are able to recursively generate transpositions of the form: 
    \begin{center}
        $(1\text { }k)=(k-1\text{ }k)^{-1}(1\text{ }k-1)(k-1\text{ }k)$
    \end{center}
    and thus the transpositions $(1\text{ }k)$, where $2 < k \leq n$ are in $\langle (1\text{ }2), s \rangle$.
    
    Finally we show that for all $1 \leq i < j \leq n$, we are able to
    create the transpositions $(i\text{ }j)$ by the product:
    \begin{center}
        $(i\text{ }j)=(1\text{ }i)^{-1}(1\text{ }j)(1\text{ }i)$.
    \end{center}
    We have shown that all possible transpositions can be generated by
    $(1\text{ }2)$ and $(1\text{ }2\text{ }\cdots\text{ }n)$ and since any permutation
    is the product of transpositions, we conclude that the set
    $S_n = \langle (1\text{ }2)$ and $(1\text{ }2\text{ }\cdots\text{ }n) \rangle $.
\end{proof}
    
    

\section*{Problem 3}
\noindent\textit{A subgroup $H$ of $S_n$ is called transitive if for all
$i,j \in \{1,2,\cdots,n\}$, there exists an element $\sigma \in H$ such that
$\sigma(i)=j$.}

\noindent\text{(a)} Show that for any group $G$ of order $n$, there exists a transitive subgroup
$H$ of $S_n$ isomorphic to $G$.

\begin{proof}
    By Caley's Theorem, there always exists $H$ a subgroup of $S_n$ which is isomorphic
    to group $G$ of finite order. It is always possible to
    define an injective homomorphism $\phi: G \to S_n$ where element $g \in G$ is sent to $gx$ for all $x \in G$.
    This subgroup is transitive because the image of $\phi$ acts transitively on the elements of $G$, i.e.
    for any two elements $g_i,g_j \in G$, there exists some $g\in G$ such that $gg_i=g_j$.
\end{proof}


\noindent\text{(b)} In the case $G=S_3$, describe a homomorphism $f:S_3 \to S_6$
with image $H$ as in part (a) by giving the image of $(1\text{ }2)$ and $(2\text{ }3)$
in cycle notation.


\noindent \subsubsection*{Solution}
We first describe a function $\phi: S_3 \to Z_6$ which assigns a label to each member of 
$S_3$. Each label is an integer between 1 and $6$.
\begin{table}[h!]
    \begin{center}
        \caption{label function $\phi$}
        \begin{tabular}{l|c}
            \textbf{$S_3$} & \textbf{$n$} \\
            \hline
            $()$ & 0 \\
            $(1\text{ }2\text{ }3)$ & 1\\
            $(2\text{ }1\text{ }3)$ & 2\\
            $(1\text{ }2)$ & 3\\
            $(1\text{ }3)$ & 4\\
            $(2\text{ }3)$ & 5\\
        \end{tabular}
    \end{center}
\end{table}

$\phi$ is an injective homomorphism because for each element of $S_n$ maps to a 
unique $n$ and $\phi(s_1s_2) = \phi(s_1)\phi(s_2)$ for all $s_1,s_2 \in S_3$ since $S_3$ is
a group under composition. Next we define a homomorphism $f: S_3 \to S_6$. $f$ is an injective homomorphism
because $f(e_{S_3}) = e_{S_6}$ since composing every element of $S_3$ with the identity then
mapping to $S_6$ is the same as mapping $S_3$ to $S_6$ then composing with the identity. The
image of $f$ acts transitively on $S_3$ since for every $g_i,g_j \in S_3$ there exists
a $g \in S_3$ such that $gg_i = g_j$.

Given the functions $\phi$ and $f$, image of $f(1\text{ }2)$ is: $(1\text{ }4)(2\text{ }5)(3\text{ }6)$. 

Given the functions $\phi$ and $f$, image of $f(2\text{ }3)$ is: $(1\text{ }6)(2\text{ }4)(3\text{ }5)$.

\section*{Problem 4}
\noindent\textit{Let $n_1,n_2,\cdots,n_k$ be positive integers
summing to $n$. Using Lagrange's Theorem, show that $n!$ is divisible
by $\prod_{i=1}^{k}n_i!$.}

\begin{proof}
    We group the integers from 1 to $n$ into $k$ partitions $H_i$, where
    $1 \leq i \leq k$. Let $A(H_i)$ be the set of all permutations
    of $H_i$. Given that $ord(H_i) = n_i$ then we have that
    $ord(A(H_i))= n_i!$ and $A(H_i) \cong S_{n_i}$.
    
    We define $H=A(H_1)\times \cdots \times A(H_k)$. The order of $H$
    is $\prod_{i=1}^{k}n_i!$. Given that each $A(H_i)$ is 
    closed under composition and each element has an inverse element, we have that
    $H$ is a subgroup of $G$, where $G=S_n$. By Lagrange's Theorem
    Theorem, since $H$ is a finite subgroup of $G$, then $|H|$ divides $|G|$. 
    Thus we conclude that $\prod_{i=1}^{k}n_i!$ divides $n!$.

\end{proof}

\section*{Problem 5}
\noindent\textit{A puzzle consists of discs labeled 1 through 99 in a circular track.
The discs can slide freely along the track (keeping their cycle order). At one point in the track, there is 
a turntable that can reverse the order of four consecutive discs. Suppose 
that initially the numbers are in order from 1 to 99 going clockwise. Is it
possible to rearrange the numbers to be in order from 1 to 99 going
counterclockwise?}

\begin{proof}
We start the puzzle at the permutation $e=()$ which is an even permutation.
The only 3 possible moves at any given moment to make progress are $t = (1\text{ }4)(2\text{ }3)$,
$r = (1\text{ }2\text{ }\cdots 99)$, and $r^{-1} = (1\text{ }2\text{ }\cdots 99)^{-1}$, all of
which are also even permutations since $t$ is the product of two transpositions and $r$ and $r^{-1}$ can both be decomposed
into 98 transpositions. We know that an even permutation composed with an even permutation
is an even permutation, and a permutation cannot be both even and odd. Given that the final
objective state of the puzzle is the permutation $ccw=(1\text{ }99)(2\text{ }98)\cdots(49\text{ }51)$, we see
that $ccw$ is an odd permutation as it is the product of 49 transpositions. Since $ccw$ cannot be reached by composition of
even permutations then we conclude that it is not possible to rearrange the discs
counterclockwise.
\end{proof}

\end{document}